Source : DGFiP, Demandes de valeurs foncières (DVF) 2024, version géolocalisée Etalab ; calculs de l'auteur. Terrains à bâtir : ventes dont toutes les parcelles sont classées « terrain à bâtir » au cadastre, sans local, de 100 à 5 000 m² (27 353 ventes). Maisons : ventes d'une seule maison sans autre logement ni local d'activité (304 799 ventes).\par Lecture : départements comptant au moins 30 ventes de terrains à bâtir ; l'Alsace-Moselle (livre foncier) est absente de DVF, Paris et les Hauts-de-Seine n'ont presque pas de terrains nus vendus. La part du terrain est le prix médian du m² de terrain multiplié par la surface médiane de parcelle des maisons vendues, rapporté au prix médian de ces maisons : un ordre de grandeur, pas une mesure.
